% ==============================================================================
\input{before.tex}\title{   Knowledge Journal   }\begin{document} \maketitle \tableofcontents
% ==============================================================================





% ==============================================================================
\clearpage\section{ Isolating CPU Cores }


\begin{enumerate}

\item isolating CPU-cores                                                                                                                                                     
    \tab "isolcpus" is the traditional, built-in method to isolate the cpu from the scheduler.                                                                                     \tab essentially allows us to prevent the OS-scheduler from assigning any process or thread to our core.                                                                  
    \tab It remains idle until we assign a program/process to it explicitly using taskset or numactl                                                                          
    \tab To use this, open /etc/default/grub, and add isolcpus=2,3 (to isolate cpu 2 and 3) to GRUB_CMDLINE_LINUX_DEFAULT string                                              
    \tab GRUB_CMDLINE_LINUX_DEFAULT="quiet splash isolcpus=2". Save the file, exit.    
    \tab update the GRUB bootloader using "sudo update-grub"                                                                                                   
    \tab reboot the system for the changes to take place. sudo reboot.                                                                                                    
    \tab also isolcpus only prevents user-tasks from being assigned there. It might still use it for hardware interrupts.                                     
    \tab to prevent this, move the hardware interrupts to somewhere else, like core 0( Make sure to do this after rebooting). 
    \tab To do this, "echo 1 $\vert$ sudo tee /proc/irq/default_smp_affinity"
    \tab now that the core is just chilling, assign your binary executable to it in the following manner                                                       
    \tab taskset -c 2 ./your-binary (Naturally, the 2 indicates the core to which you wanna write to this)
                                                                                                                                                                             
                                                                                                                                                                              \item undoing the cpu-core isolation                                                                                                                          
    \tab sudo nano /etc/default/grub                                                                                                                              
    \tab remove the isolcpus=2 fom the GRUB_CMDLINE_LINUX_DEFAULT string                                                                                           
    \tab sudo update-grub                                                                                                                                             
    \tab sudo reboot                                                                                                                                                           
\end{enumerate}




% ==============================================================================
\end{document}
% ==============================================================================
